\section{Resultados e Síntese}

Esta seção consolida os achados dos 37 estudos incluídos. A síntese segue uma abordagem temática em modo híbrido: seis temas foram derivados indutivamente dos campos estruturados das fichas de extração (artefatos, papéis, execução, validação, portabilidade, riscos, natureza e evidência) e, em seguida, confrontados com a taxonomia de seis dimensões proposta na seção anterior, de modo a testá-la contra a evidência em vez de assumi-la. Os temas respondem às perguntas RQ1 a RQ4; a RQ5, sobre lacunas, é tratada na Agenda de Pesquisa. A força da evidência de cada estudo é ponderada por uma avaliação de qualidade por categoria de fonte, descrita no Método e resumida adiante (Tabela~\ref{tab:qualidade}).

\noindent\textbf{Panorama quantitativo.} A Tabela~\ref{tab:quant} resume a composição do corpus. A categoria framework domina quase por completo (34 dos 37 estudos), o que sustenta empiricamente a tese de que a unidade de trabalho se deslocou do prompt isolado para o processo estruturado. A avaliação empírica é reivindicada por 26 estudos e 22 disponibilizam repositório público, mas dez são preprints arXiv ainda não revisados por pares, o que exige leitura qualificada da força da evidência (Figura~\ref{fig:evidence}). A distribuição temporal mostra aceleração consistente do campo, de um único estudo em 2022 para onze em 2026, com a maior parte do corpus concentrada a partir de 2024 (Figura~\ref{fig:year}). Os venues são dispersos, com a ACM Transactions on Software Engineering and Methodology como canal recorrente (cinco estudos) seguido de uma cauda longa de conferências e periódicos com um a dois estudos cada, o que é coerente com um campo emergente ainda sem foro consolidado.

\begin{table}[h!]
\centering
\caption{Composição do corpus de 37 estudos incluídos. Contagens extraídas da matriz de extração.}
\label{tab:quant}
\small
\begin{tabularx}{\textwidth}{@{}p{0.50\textwidth}r Y@{}}
\toprule
\textbf{Aspecto} & \textbf{N} & \textbf{Observação} \\
\midrule
Total de estudos incluídos & 37 & Após triagem em duas passadas e três rodadas de snowballing \\
Natureza: framework & 34 & Predominância quase total \\
Natureza: arquitetura / plataforma / método & 1 / 1 / 1 & JulesWhite2023, YuCheng2023, JSauvola2024 \\
Com avaliação empírica & 26 & Algum grau de validação reivindicado \\
Com repositório público & 22 & Rastreabilidade de código em mais da metade \\
Preprint arXiv & 10 & Literatura ainda não revisada por pares \\
Survey ou revisão & 2 & AhmedEHassan2024, JundaHe2025 \\
Estudo de caso & 5 & Ulfsnes2024, ZeeshanRasheed2024, Baron2025, Khan2026, Macedo2026 \\
Distribuição por ano & \multicolumn{2}{l}{2022:1, 2023:4, 2024:10, 2025:11, 2026:11} \\
Origem acadêmica ou técnica / industrial & 33 / 4 & Quatro estudos de indústria ou prática profissional \\
Confiança de extração alta & 35 & Média em Dam2025 e JundaHe2025 \\
\bottomrule
\end{tabularx}
\end{table}

\begin{figure}[h!]
\centering
\includegraphics[width=0.70\textwidth]{studies_by_year_pt.pdf}
\caption{Estudos por ano de publicação. O campo acelera de um único estudo em 2022 para um patamar de dez a onze por ano a partir de 2024.}
\label{fig:year}
\end{figure}

\begin{figure}[h!]
\centering
\includegraphics[width=0.80\textwidth]{evidence_profile_pt.pdf}
\caption{Perfil de evidência do corpus. A maioria dos estudos reivindica avaliação empírica e disponibiliza repositório público, mas dez são preprints arXiv ainda não revisados por pares.}
\label{fig:evidence}
\end{figure}

\noindent\textbf{Avaliação de qualidade por categoria de fonte.} A Tabela~\ref{tab:qualidade} resume a avaliação de qualidade descrita no Método, com o checklist de quatro itens aplicado aos 37 estudos. Onze estudos (30\%) ficam na faixa de qualidade alta, dezessete (46\%) na média e nove (24\%) na baixa. A maior parte do corpus vem de venue revisado por pares (25 estudos, 68\%), contra dez preprints arXiv (27\%) e dois registros técnicos não indexados (5\%). No nível dos itens, 26 estudos (70\%) reportam avaliação empírica e 22 (59\%) disponibilizam repositório público, mas apenas cinco (14\%) trazem validação aplicada em estudo de caso ou projeto real, o que delimita o quanto a evidência atual sustenta uso em produção. A faixa baixa concentra-se em preprints e registros técnicos sem avaliação empírica registrada; seus achados são lidos com cautela na síntese e nenhum tema se apoia exclusivamente neles. A avaliação é conservadora por construção, pois deriva dos sinais registrados na extração e não de uma reanálise de texto completo estudo a estudo.

\begin{table}[h!]
\centering
\caption{Avaliação de qualidade por categoria de fonte dos 37 estudos. Itens derivados de sinais observáveis na matriz de extração; QA é a soma de QA1 a QA4 (0 a 4).}
\label{tab:qualidade}
\small
\begin{tabularx}{\textwidth}{@{}p{0.46\textwidth}r Y@{}}
\toprule
\textbf{Categoria} & \textbf{N} & \textbf{\%} \\
\midrule
Qualidade alta (QA $\geq$ 3) & 11 & 30\% \\
Qualidade média (QA = 2) & 17 & 46\% \\
Qualidade baixa (QA $\leq$ 1) & 9 & 24\% \\
\midrule
Venue revisado por pares (QA1) & 25 & 68\% \\
Preprint arXiv & 10 & 27\% \\
Registro técnico não indexado & 2 & 5\% \\
\midrule
Avaliação empírica reportada (QA2) & 26 & 70\% \\
Validação aplicada em estudo de caso (QA3) & 5 & 14\% \\
Artefato público para reprodutibilidade (QA4) & 22 & 59\% \\
\bottomrule
\end{tabularx}
\end{table}

\noindent\textbf{Seis temas indutivos.} A análise fez emergir seis temas que, juntos, cobrem os 37 estudos (cobertura de 100\%, acima do alvo de 80\%). Cada estudo pode pertencer a mais de um tema; a ordem reflete frequência decrescente.

\noindent\textbf{T1: orquestração multiagente e divisão de papéis.} Este é o tema mais populoso e o eixo arquitetônico dominante do corpus, respondendo diretamente à RQ1. A convergência central é que a estruturação do trabalho por papéis distintos, e não a escolha do modelo de linguagem, é o que diferencia os frameworks. ChatDev codifica uma cadeia de conversas com fases de design, codificação, teste e revisão, além de um mecanismo de redução comunicativa de alucinação \cite{ChenQian2023}. MetaGPT internaliza procedimentos operacionais padrão de engenharia de software, atribuindo a cada papel artefatos específicos: o gerente de produto gera requisitos, o arquiteto gera o design e as interfaces, o desenvolvedor implementa \cite{SiruiHong2023}. CodePori escala esse arranjo para agentes de gerência, arquitetura, estrutura, desenvolvimento, verificação e finalização \cite{ZeeshanRasheed2024}, e AutoDev e AutoCodeRover levam a coordenação de agentes a tarefas reais de melhoria de programas \cite{MicheleTufano2024,YuntongZhang2024}. As revisões de He et al. e a visão de Engenharia de Software 3.0 de Hassan et al. confirmam que a literatura já organiza esses sistemas por etapas do ciclo de vida \cite{JundaHe2025,AhmedEHassan2024}. O consenso, porém, é parcial: há tensão clara quanto ao grau de autonomia. Sapkota et al. contrastam explicitamente codificação por intuição e codificação agêntica \cite{RanjanSapkota2025}, e enquanto CodePori e MetaGPT se inclinam ao desenvolvimento autônomo de ponta a ponta, outros desenhos preservam pontos de decisão humana. A confiabilidade e o custo de coordenação aparecem como riscos recorrentes neste tema \cite{Dam2025,Alsegier2026,Merchant2025,Paduraru2026,SKaruppuchamy2026,Paiva2026,Akbar2025,LiyiCai2025,Anon2026,BESLEAGA2026}.

\noindent\textbf{T2: especificação e contrato como artefato central.} Este tema responde RQ1 e RQ2 e materializa a tese de que a especificação se torna a unidade de coordenação. O desenvolvimento orientado por especificação formaliza a passagem do código ao contrato, tratando a especificação como artefato versionável \cite{DeepakBabuPiskala2026}. AgileGen usa critérios de aceitação como contrato executável e mantém memória e iteração até a validação \cite{SaiZhang2025}, e MetaGPT já fazia dos requisitos estruturados a fonte de verdade entre papéis \cite{SiruiHong2023}. Há, contudo, uma contradição produtiva quanto à direção do fluxo: a maioria assume \textit{spec-first}, em que a especificação precede e gera o código, ao passo que Reversa inverte o fluxo, recuperando especificações operacionais a partir de código legado para alimentar os agentes \cite{Macedo2026}. Essa oposição entre especificar antes e recuperar a especificação do legado é um eixo de diferenciação relevante. A engenharia de requisitos consciente de risco liga a especificação à governança \cite{Watfa2026}, e o desalinhamento entre especificação e código é discutido como risco em vários estudos \cite{DeepakBabuPiskala2026,Paduraru2026,Watfa2026,YuCheng2023,Gupta2022,MitulModi2024,Khan2026,Gadamsetty2025}.

\noindent\textbf{T3: engenharia de contexto e recuperação de conhecimento do repositório.} Este tema responde RQ2 com foco na dimensão de contexto. A engenharia de contexto é tratada como disciplina própria para agentes em software de código aberto \cite{SeyedmoeinMohsenimofidi2025}. AutoCodeRover demonstra que a recuperação estrutural de contexto no código, e não apenas a busca textual por embeddings, melhora a geração de correções \cite{YuntongZhang2024}, enquanto a sincronização de artefatos multimodais busca manter coerência entre artefatos heterogêneos ao longo do ciclo \cite{Paiva2026}. O software autoevolutivo leva o tema ao extremo \cite{LiyiCai2025}, mas entra em tensão com a ênfase em rastreabilidade e contrato dos temas T2 e T4: quanto mais o software se modifica de forma autônoma, mais difícil é ancorar cada mudança em uma especificação auditável \cite{Macedo2026,AhmedEHassan2024,Chechik2026}. O risco associado ao contexto é o mais frequente do corpus inteiro, o que sugere que a janela e a fidelidade de contexto são o gargalo prático transversal do campo.

\noindent\textbf{T4: validação, rastreabilidade, asseguramento e governança.} Este tema responde RQ2, na dimensão de validação, e RQ4, nos riscos de governança e rastreabilidade. Quase todo o corpus declara algum mecanismo de validação, mas aqui se agrupam os estudos para os quais a validação é o ponto central, e não um apêndice. O asseguramento baseado em traços articula contratos, testes e governança \cite{Paduraru2026}; a colaboração verificável entre assistentes agênticos registra atos do fluxo de forma auditável \cite{Merchant2025}; e a engenharia de linha de produto é importada para governar a variabilidade de agência \cite{Alsegier2026}. A rastreabilidade e a auditoria de traços aparecem como requisito recorrente para confiar em agentes em projetos reais \cite{Paduraru2026,Merchant2025,Macedo2026,SKaruppuchamy2026,Watfa2026,Chechik2026,Anon2026}. Há, porém, uma contradição de custo: de um lado, cerimônia pesada de governança \cite{Merchant2025,Paduraru2026}; de outro, agilidade leve voltada a startups e a fluxos enxutos \cite{Khan2026,RanjanSapkota2025,JulesWhite2023}. O custo de coordenação é nomeado como risco em boa parte deste grupo \cite{MitulModi2024,Ulfsnes2024,Akbar2025,Baron2025,KRRaghi2024,Gupta2022,KumarAle2024}, o que materializa a tensão entre governança e produtividade.

\noindent\textbf{T5: colaboração humano-IA, adoção e confiança (emergente).} Estes estudos têm como foco primário não a arquitetura técnica, mas a integração da IA no trabalho humano: fluxos de colaboração, fatores de adoção organizacional e construção de confiança. Um framework de colaboração e adaptação humano-IA modela fatores de fluxo de trabalho e estratégias organizacionais \cite{DanielRusso2024}; evidências de campo descrevem como equipes reais reorganizam colaboração e fluxo \cite{Ulfsnes2024,Akbar2025}; e a confiança é tratada como precondição de adoção \cite{Baron2025}, com cenários futuros de operação assistida por IA também documentados \cite{JSauvola2024}. O achado relevante é que este tema não cabe limpo em nenhuma das seis dimensões da taxonomia, todas técnico-arquiteturais: a dimensão humana, organizacional e de confiança atravessa o corpus mas não tem casa na taxonomia atual. Seis estudos (16\% do corpus) têm este foco primário, material o bastante para recomendar uma extensão da taxonomia.

\noindent\textbf{T6: segurança, robustez e red-teaming (emergente).} Estes estudos focam a segurança do código e do comportamento agêntico. ASTRA propõe \textit{red-teaming} espaço-temporal autônomo para assistentes de software com IA, tratando segurança como atividade ativa e não como propriedade assumida \cite{XiangzheXu2025}; a engenharia de requisitos consciente de risco e a discussão das implicações de segurança da codificação agêntica reforçam o tema \cite{Watfa2026,RanjanSapkota2025,Merchant2025,BESLEAGA2026}, e a redução comunicativa de alucinação de ChatDev tangencia o ponto pelo lado da confiabilidade \cite{ChenQian2023}. A segurança está entre os riscos mais citados, nomeada em 24 dos 37 estudos, atrás de contexto, confiabilidade e custo de coordenação (Figura~\ref{fig:risk}). Assim como T5, este tema só se conecta parcialmente às seis dimensões: segurança aparece quase sempre como risco a mitigar, raramente como dimensão de projeto positiva embutida na arquitetura. ASTRA é a exceção que faz da segurança o objeto central, o que sugere tratá-la como critério de avaliação transversal mais do que como uma das dimensões que classificam frameworks.

\begin{figure}[h!]
\centering
\includegraphics[width=0.80\textwidth]{risk_frequency_pt.pdf}
\caption{Frequência dos riscos reportados no corpus. Contexto, confiabilidade e custo de coordenação lideram; segurança, embora proeminente, fica abaixo deles, e a marcação explícita de aprisionamento é rara.}
\label{fig:risk}
\end{figure}

\noindent\textbf{Confronto da taxonomia com a evidência.} A Tabela~\ref{tab:tema-dim} cruza cada tema indutivo com as seis dimensões da taxonomia, indicando a força da relação, codificada pelo autor a partir das fichas de extração. A leitura é direta. Cinco das seis dimensões recebem suporte forte de pelo menos um tema: papéis e execução ancoram-se em T1; especificação e validação, em T2 e T4; contexto, em T3. Nessa parte, a taxonomia se sustenta no confronto com a evidência. A Figura~\ref{fig:heatmap} apresenta o mesmo mapeamento como um mapa de calor. A exceção é portabilidade, que não recebe relação forte de nenhum tema e é ausente em dois. Nos dados, portabilidade aparece quase sempre pelo lado negativo, como risco de aprisionamento tecnológico e de dependência de IDE, modelo comercial ou plataforma. O campo de portabilidade das fichas registra dependência de IDE ou ambiente em praticamente todos os 37 estudos, e o aprisionamento só é nomeado explicitamente como risco em um deles \cite{SeyedmoeinMohsenimofidi2025}. Portabilidade é, portanto, a dimensão mais fraca da taxonomia quando confrontada com a evidência, e deve ser lida como dimensão predominantemente diagnóstica, e não como diferenciador positivo entre frameworks. Os temas T5 e T6, por sua vez, não têm casa clara nas seis dimensões: o eixo humano-organizacional e o eixo de segurança emergem dos dados como recomendações concretas de extensão, e a regra do modo híbrido foi respeitada ao não forçá-los nas dimensões existentes.

\begin{table}[h!]
\centering
\caption{Mapeamento híbrido: temas indutivos contra as seis dimensões da taxonomia. Legenda: forte (eixo do tema), média (presença relevante), fraca (toque marginal), ausente (sem relação significativa).}
\label{tab:tema-dim}
\small
\begin{tabularx}{\textwidth}{@{}p{3.5cm}YYYYYY@{}}
\toprule
\textbf{Tema} & \textbf{Espec.} & \textbf{Contexto} & \textbf{Papéis} & \textbf{Execução} & \textbf{Validação} & \textbf{Portab.} \\
\midrule
T1 Orquestração multiagente & média & média & forte & forte & média & fraca \\
T2 Especificação e contrato & forte & média & fraca & média & forte & fraca \\
T3 Engenharia de contexto & fraca & forte & fraca & média & fraca & fraca \\
T4 Validação e governança & média & fraca & fraca & média & forte & fraca \\
T5 Colaboração humano-IA & fraca & fraca & fraca & fraca & fraca & ausente \\
T6 Segurança e red-teaming & fraca & fraca & fraca & média & média & ausente \\
\bottomrule
\end{tabularx}
\end{table}

\begin{figure}[h!]
\centering
\includegraphics[width=0.92\textwidth]{theme_dimension_pt.pdf}
\caption{Mapeamento híbrido dos seis temas indutivos contra as seis dimensões da taxonomia. Cinco dimensões são ancoradas por pelo menos um tema; portabilidade permanece fraca em toda a linha, e os temas T5 e T6 ficam parcialmente fora da taxonomia.}
\label{fig:heatmap}
\end{figure}

\noindent\textbf{Cobertura das perguntas de pesquisa.} A Tabela~\ref{tab:tema-rq} cruza os temas com as perguntas de pesquisa. RQ1 e RQ2 são respondidas pela maioria dos temas; RQ3 e RQ4 atravessam o corpus inteiro, e a RQ5 é meta-pergunta respondida transversalmente na Agenda de Pesquisa, e não por marcação estudo a estudo, razão pela qual as fichas individuais marcam apenas RQ1 a RQ4.

\begin{table}[h!]
\centering
\caption{Temas indutivos contra as perguntas de pesquisa.}
\label{tab:tema-rq}
\small
\begin{tabularx}{\textwidth}{@{}p{3.5cm}YYYYY@{}}
\toprule
\textbf{Tema} & \textbf{RQ1} & \textbf{RQ2} & \textbf{RQ3} & \textbf{RQ4} & \textbf{RQ5} \\
\midrule
T1 Orquestração multiagente & sim & sim & sim & sim & via síntese \\
T2 Especificação e contrato & sim & sim & parcial & sim & via síntese \\
T3 Engenharia de contexto & sim & sim & parcial & sim & via síntese \\
T4 Validação e governança & parcial & sim & sim & sim & via síntese \\
T5 Colaboração humano-IA & parcial & sim & sim & sim & via síntese \\
T6 Segurança e red-teaming & não & parcial & sim & sim & via síntese \\
\bottomrule
\end{tabularx}
\end{table}
